\documentclass[10pt,a4paper]{amsart}
\usepackage[margin=19mm]{geometry}
\usepackage{amsmath,amssymb,mathtools}
\usepackage[scr=rsfs,cal=euler]{mathalfa}
\usepackage{xfrac}
\usepackage[unicode,hidelinks,bookmarks=false]{hyperref}
\newcommand{\ie}{i.e.\ }
\newcommand{\cf}{cf.\ }
\newcommand{\resp}{respectively\ }
\newcommand{\Subsection}{\subsection}
\providecommand{\nobreakdash}{\nobreak\hbox{-}}
\setlength{\emergencystretch}{2em}
\multlinegap0pt
\usepackage{amssymb}
\usepackage{tikz}
\usepackage{enumerate}
\usepackage{xcolor}
\newtheorem{theorem}{Theorem}
  \newcommand{\Id}{\ensuremath{\mathop{\rm Id}\nolimits}}
  \newcommand{\R}{\mathbb{R}} % reals
  \DeclareMathOperator{\Tr}{Tr}
  \DeclareMathOperator{\cb}{cb}
  \newcommand{\ind}[1]{\mathbf{#1}}
  \newcommand{\norm}[1]{\|#1\|}
\title{Optimal inequalities for completely bounded polynomials\\ and the limitations of quantum query algorithms}
\author{Francisco Escudero Guti\'errez, Miquel Saucedo, Carlos Palazuelos}
\date{Source: arXiv:2609.05201v1 (CC BY 4.0)}
\begin{document}
\maketitle

\noindent\emph{Source and licence.} Optimal inequalities for completely bounded polynomials\\ and the limitations of quantum query algorithms. Version arXiv:2609.05201v1 (CC BY 4.0), distributed by arXiv under the Creative Commons Attribution 4.0 International licence, http://creativecommons.org/licenses/by/4.0/, as recorded in the arXiv licence metadata for this exact version and shown on the abstract page https://arxiv.org/abs/2609.05201v1. Full licence notice: \texttt{licenses/arxiv-2609.05201-CC-BY-4.0.txt}.\par\smallskip

\noindent\textit{Editorial scope.} Counted unit: the article's theorem theo:optimalcbBHBlockMult (source0.tex lines 1091--1098). The article's complete original proof of it is delivered with it (source0.tex lines 1100--1173, one \texttt{proof} environment of 74 lines). No mathematics is written, repaired or re-derived: the statement and the proof are the article's own lines, in the article's own notation, and no retained line is substituted or re-flowed.  No further source-local context is needed: every cross-reference of every retained line resolves inside the two delivered spans.  Nothing was omitted from any retained span, so the package carries no omission record.  No retained line contains a citation, so the wrapper carries no bibliography and no bibliography entry.  No retained line contains a figure environment, an image-inclusion command or drawing code, so no image is delivered and the package declares no asset dependency.  Editorial formatting only, in addition: an AMS \texttt{amsart} preamble that restates the article's own declarations and macros used by the retained lines, namely the environments \texttt{theorem} on separate counters, together with the standard packages those lines need.  The retained environments print in this document as Theorem 1 in order of appearance, whereas the article prints the same items with the numbering of its own full document.  The complete native source of the article as distributed in the arXiv e-print for this exact version is bundled unchanged as \texttt{sources/209421/source0.tex}.
\begin{theorem}\label{theo:optimalcbBHBlockMult}

Let $p:\{-1,1\}^n\times\dots\times\{-1,1\}^n\to\R$ be a block-multilinear polynomial of degree $t$. Then, for all $s\in [t]$, it is satisfied that  
\begin{equation}
\label{eq:optimalcbBHBlockMult}
    \norm{p}_{\cb}\geq   \sum_{i\in [n]_0} \left(\sum_{\ind  i\in [n]^t_0:i_s=i} |p_{\ind i}|^2\right)^\frac12. 
\end{equation}  
\end{theorem}

\begin{proof}
For a fixed $s\in [t]$, we write $\ind i=({\ind{i}_L},i_s,{\ind{i}_R})$ with ${\ind{i}_L}=(i_1,\dots, i_{s-1})$ and ${\ind{i}_R}=(i_{s+1},\dots, i_t)$.  We also write 
$$X_1(i_1)\cdots X_{s-1}(i_{s-1})=X_{L}({\ind{i}_L}) \mbox{ and }X_{s+1}(i_{s+1})\cdots X_d(i_{d})=X_R({\ind{i}_R}).$$
It may be that ${\ind{i}_L}$ or ${\ind{i}_R}$ are empty. Then $X_L$ or $X_R$ are just the identity. We will carry out the construction of the unitaries and unit vectors in the space $\mathcal{H}=\bigotimes _{l=1}^t\R^{n+1}.$ We divide the proof in 3 steps.

{\bf Step 1.} First, we define $e$ and $X_l$ for $l\geq s+1.$ 
To do so, consider the  orthonormal basis $$(e_{\ind i}:=e_{i_{1}} \otimes \cdots \otimes e_{i_{t}})_{\ind i\in[n]_0^t},$$ and set $e=e_0\otimes \cdots\otimes e_0$. Then, $X_{k}(i)$ acts only on the factor of $\mathcal{H}$ with $l=k$ (being the identity on the rest of the factors) and satisfies $X_k(i)e_0=e_{i}$.   
From the definition, we see that $(X_{R}({\ind{i}_R})e)_{{{\ind{i}_R}}}$ is an orthonormal set. 
    

 Set $$v_{{\ind{i}_L},i_s}=\sum_{{\ind{i}_R}} p_{({\ind{i}_L},i_s,{\ind{i}_R})} X_R({\ind{i}_R}) e,$$ and observe that, by orthonormality 
\begin{equation}
\label{eq:orthonormalityv}\norm{v_{{\ind{i}_L},s}}_{\ell_2}= \left(\sum_{{\ind{i}_R}}  |p_{({\ind{i}_L},i_s,{\ind{i}_R})}|^2\right)^\frac12.\end{equation}

  {\bf Step 2.} Second, we claim that
\begin{equation}
\label{eq:step1}\norm{p}_{\cb}\geq \sum_{i_s=0}^n \sup _{\substack{\mathcal{F}=(f_{{\ind{i}_L}})_{{\ind{i}_L}}\\
\mathcal{F}\text{ orthonormal}\\|\mathcal{F}|=(n+1)^{s-1}}} \sum_{{\ind{i}_L}} \langle f_{{\ind{i}_L}}, v_{{\ind{i}_L},i_s}\rangle.\end{equation}


To see this, fix $n+1$ orthonormal sets $\mathcal{F}_0, \dots, \mathcal{F}_{n}$ of size $(n+1)^{s-1}$, and set $\mathcal{F}_m=(f^m_{{\ind{i}_L}})_{{\ind{i}_L}}$ for every $m=0,\ldots, n$. A simple computation shows that for any vector $f$ and any matrices  $X_1(i_1),\dots, X_s(i_s)$, we have that 
    \begin{align*}
        \langle f, p(X) e\rangle &=\sum_{i_s=0}^ n \left( \sum_{{\ind{i}_L}} \langle f, X_L({\ind{i}_L}) X_s(i_s) v_{{\ind{i}_L},i_s}\rangle\right) \\
         &=\sum_{i_s=0}^ n \left( \sum_{{\ind{i}_L}} \langle X_s(i_s) ^{\mathsf T} X_{L}({\ind{i}_L})^{\mathsf T} f,   v_{{\ind{i}_L},i_s}\rangle\right).
    \end{align*}


Now, since all the orthonormal sets have the same size, for any $m$ we can find an orthogonal matrix $X_s(m)$ such that $$X_s(m)(f^m_{\ind{i}_L})=f^0_{\ind{i}_L}\hspace{0.4 cm}\text{for every $\ind{i}_L$},$$ with  $X_s(0)=\Id$. 



On the other hand, we can easily find orthogonal matrices $X_1(i_1),\ldots, X_{s-1}(i_{s-1})$, with $X_k(0)=\Id$ for $k=1,\ldots, s-1$ such that $$X_L ({i}_L)f^0_{{i}_L}=f_{0},$$ where $f_{0}=f^0_{(0,\dots,0)}$.  Indeed, in order to define these matrices, consider, for every $m=0,\ldots, n$ and $k=1,\ldots, s-1$, the orthonormal sets $$(f^0_{i_1,\ldots , i_{k-1}, m,0,\ldots, 0})_{i_1,\ldots, i_{k-1}=1}^n\hspace{0.3 cm} \text{and}\hspace{0.3 cm}(f^0_{i_1,\ldots , i_{k-1}, 0,0,\ldots, 0})_{i_1,\ldots, i_{k-1}=1}^n.$$ Since they have the same size, we can always find orthogonal matrices, $X_k(m)$ with $X_k(0)=\Id$ for every $k$ and satisfying $$X_k(m)(f^0_{i_1,\ldots , i_{k-1}, m,0,\ldots, 0})=f^0_{i_1,\ldots , i_{k-1}, 0,0,\ldots, 0}.$$ It is clear that these matrices satisfy the desired property.



    Then,

    $$\langle f_{0}, p(X) e\rangle =\sum_{i_s=0}^ n \left( \sum_{{\ind{i}_L}} \langle X_s(i_s) ^{\mathsf T} f^0_{{\ind{i}_L}},   v_{{\ind{i}_L},i_s}\rangle\right)=\sum_{i_s=0}^ n \left( \sum_{{\ind{i}_L}} \langle  f^{i_s}_{{\ind{i}_L}},   v_{{\ind{i}_L},i_s}\rangle\right).$$ 
    As a consequence,
 $$\norm{p}_{\cb}\geq \sum_{i_s=0}^n \sup _{\substack{\mathcal{F}_{i_s}=(f^{i_s}_{{\ind{i}_L}})_{{\ind{i}_L}}\\
  \mathcal{F}_{i_s}\text{ orthonormal}\\|\mathcal{F}_{i_s}|=(n+1)^{s-1}}} \sum_{{\ind{i}_L}} \langle f^{i_s}_{{\ind{i}_L}}, v_{{\ind{i}_L},i_s}\rangle,$$ which is the same as \eqref{eq:step1}.

   

 {\bf Step 3.} Third, we show that for any vectors $(w_{{\ind{i}_L}})_{{\ind{i}_L}}$

\begin{equation}
    \label{eq:step2}
 \sup _{\substack{\mathcal{F}=(f_{{\ind{i}_L}})_{{\ind{i}_L}}\\
\mathcal{F}\text{ orthonormal}\\|\mathcal{F}|=(n+1)^{s-1}}} \sum_{{\ind{i}_L}} \langle f_{{\ind{i}_L}}, w_{{\ind{i}_L}}\rangle \geq \left(\sum_{{\ind{i}_L}} \|w_{{\ind{i}_L}}\|^2\right)^\frac12.\end{equation}


To see this, for any given $\mathcal{F}$, let $M_{\mathcal{F}}$ be the linear operator on $\mathcal{H}$ defined by $M_{\mathcal{F}}(f_{{\ind{i}_L}})={w_{{\ind{i}_L}}}$ and $M_{\mathcal{F}}(v)=0$ for $v\in \langle \mathcal{F} \rangle ^\perp$. Then,
$$ \sum_{{\ind{i}_L}} \langle f_{{\ind{i}_L}}, w_{{\ind{i}_L}}\rangle =\Tr(M_{\mathcal{F}}).$$
We remark that for any given orthonormal set $\mathcal{F}_0$ we have
$$\{ M_{\mathcal{F}}: \mathcal{F} \mbox{ orthonormal}\}= \{ M_{\mathcal{F}_0} O: O \mbox{ orthogonal}\}.$$ Indeed, let $\mathcal{F}$ be another orthonormal set of $(n+1)^{s-1}$ elements. Then there is an orthogonal matrix $O$ such that $O \mathcal{F}= {\mathcal{F}_0}$, which implies that $M_{\mathcal{F}}=M_{\mathcal{F}_0} O.$ For the reverse inclusion, given $O$ orthogonal, define $\mathcal{F}=O^{-1}\mathcal{F}_0.$




As a consequence, $$ \sup _{\substack{\mathcal{F}=(f_{{\ind{i}_L}})_{{\ind{i}_L}}\\\mathcal{F}\text{ orthonormal}\\|\mathcal{F}|=(n+1)^{s-1}}} \sum_{{\ind{i}_L}} \langle f_{{\ind{i}_L}}, w_{{\ind{i}_L}}\rangle= \sup_{O \mbox{ orthogonal}} \Tr(M_{\mathcal{F}_0}O)=\|M_{\mathcal{F}_0}\|_{S_1}.$$ 
To obtain \eqref{eq:step2} we use the fact that 



    \begin{equation}
    \label{eq:stepschatten}
        \|M_{\mathcal{F}_0}\|_{S_1} \geq \|M_{\mathcal{F}_0}\|_{S_2} =  \Tr( M_{\mathcal{F}_0}^{\mathsf T} M_{\mathcal{F}_0})^\frac12= \left(\sum_{{\ind{i}_L}} \|w_{{\ind{i}_L}}\|^2\right)^\frac12.
    \end{equation} Here, $\|M_{\mathcal{F}_0}\|_{S_1}$ stands for the trace norm of the matrix $M_{\mathcal{F}_0}$, which is known to be larger than or equal to its Hilbert-Schmidt norm.
    
Finally, the result follows by combining \eqref{eq:orthonormalityv}, \eqref{eq:step1} and \eqref{eq:step2}.


\end{proof}

\end{document}
